\section{Transform Derivations and Resumable Traversal}
\label{sec:legacy-transform}
The frame and normalization convention is defined in
Section~\ref{sec:signal-model}. This appendix records the conventional
transform mathematics and the original state transition used by the reusable
FFT classes; the complete analysis schedule is in Section~\ref{sec:one-hop}.

\subsection{Resumable complex transform}
For $N=2^p$, bit reversal places the input in the order required by an iterative
radix-2 decimation-in-time transform. A stage of span $s=2,4,\ldots,N$ visits
groups $g=0,s,2s,\ldots,N-s$ and pairs $a=0,\ldots,s/2-1$. One step computes
\begin{align}
 e&=c[g+a], & o&=W_N^{aN/s}c[g+a+s/2],\\
 c[g+a]&\leftarrow e+o, & c[g+a+s/2]&\leftarrow e-o,
 \label{eq:butterfly}
\end{align}
where $W_N=\exp(-j2\pi/N)$. Advancing $a$, then $g$, then $s$ defines a serial
traversal of the conventional FFT graph. Retaining these three indices is
sufficient to suspend and resume that traversal. Each stage has $N/2$
butterflies, so the total count is
\begin{equation}
 B_c(N)=\frac{N}{2}\log_2 N.
 \label{eq:countcomplex}
\end{equation}
The coefficient array, twiddle table, and bit-reversal table each occupy
$O(N)$ storage; the traversal state occupies $O(1)$ storage. Buffer preparation
copies and permutes the entire frame. Resumability applies to the subsequent
butterfly traversal, not automatically to that preparation.

\begin{proposition}[Arithmetic-order preservation]
For a fixed buffered frame and transform configuration, any partition of the
serial butterfly traversal into finite batches yields the same completed
coefficient array as uninterrupted traversal, provided every butterfly executes
exactly once in the original order and no external code mutates the state.
\end{proposition}
\begin{proof}
After buffering, the two traversals have identical coefficients and indices.
If they agree before a butterfly, Equation~\eqref{eq:butterfly} reads the same
operands and performs the same updates; the index transition also agrees.
Induction over $B_c$ establishes equality. A pause changes no state. In floating
point, bitwise equality additionally assumes the same arithmetic instructions,
rounding environment, and compiler transformations. The experiments test
bitwise equality within a single executable rather than asserting it across
different compilers or FFT libraries.
\end{proof}

\begin{lstlisting}[float=tbp,language=C++,caption={Simplified state transition corresponding to the original complex FFT. Tables and storage are prepared separately.},label={lst:step}]
void step() {
    if (span > N) return;
    const size_t half = span / 2;
    const auto even = c[group + pair];
    const auto odd = multiply(c[group + pair + half],
                              twiddle[pair * (N / span)]);
    c[group + pair] = even + odd;
    c[group + pair + half] = even - odd;
    if (++pair == half) {
        pair = 0;
        group += span;
        if (group == N) { group = 0; span *= 2; }
    }
}
\end{lstlisting}

\subsection{Real-input packing and reconstruction}
For even $N$, define $M=N/2$ and pack
$z[m]=u[2m]+j u[2m+1]$. Let $Z$ be its $M$-point complex transform.
For $1\le k<M$, with
$A=Z[k]$ and $D=\overline{Z[M-k]}$, reconstruction is
\begin{equation}
 X[k]=\frac{1}{2}\bigl(A+D-jW_N^k(A-D)\bigr),\qquad
 X[N-k]=\overline{X[k]}.
 \label{eq:real}
\end{equation}
The endpoints are
$X[0]=\Re Z[0]+\Im Z[0]$ and
$X[M]=\Re Z[0]-\Im Z[0]$.
This follows by separating the transforms of the even and odd real samples;
it is a standard consequence of conjugate symmetry, not a new real FFT.

The resumable inner transform performs
\begin{equation}
 B_r(N)=B_c(N/2)=\frac{N}{4}\log_2(N/2)
 \label{eq:countreal}
\end{equation}
butterflies. The original real-FFT class materializes all $N$ output bins, including conjugate
copies, and reconstructs them in one $O(N)$ loop when the last inner butterfly
finishes. Thus a real-FFT \code{step()} is not uniformly a single butterfly:
the completing step also performs reconstruction. The operational distinction
is essential to the timing model in Appendix~\ref{sec:cost}.

The inverse class uses
$\operatorname{IFFT}(X)=\overline{\operatorname{FFT}(\overline X)}/N$.
In the evaluated revision it treats each output normalization as an additional
resumable step, giving $B_c(N)+N$ steps after buffering. This differs from the
real forward transform's bulk reconstruction. Neither production visualization
module uses the inverse transform, and the timing campaign concerns only the
real forward path.
